\subsection{System Architecture Overview}
\label{sec:system-architecture}

The contagion simulation extends the existing ChatDev/DevAll platform through a modular, composable architecture that integrates seamlessly with the PixiJS-based spatial visualization layer. This section describes the technical implementation, emphasizing how the simulation engine integrates with existing rendering and agent management systems without requiring backend modifications.

\subsubsection{PixiJS-Based SpatialCanvas Environment}
\label{sec:spatialcanvas}

The \texttt{SpatialCanvas.vue} component serves as the primary rendering surface for the spatial environment, built on PixiJS 8.x for hardware-accelerated 2D graphics rendering. The canvas architecture follows a layered container model:

\begin{itemize}
    \item \textbf{Floor Layer} (\texttt{floorContainer}): Renders background tiles with optional contamination overlays, managed by \texttt{useFloorManager.js}
    \item \textbf{Grid Layer} (\texttt{gridGraphics}): Displays the 40-pixel grid overlay for spatial alignment and obstacle placement
    \item \textbf{Obstacle Layer} (\texttt{obstacleContainer}): Renders walls, furniture, and environmental barriers with collision detection
    \item \textbf{Agent Layer} (\texttt{agentContainer}): Contains agent sprites with status glows, emote bubbles, and animation state machines
    \item \textbf{Trail Layer} (\texttt{trailGraphics}): Renders particle trails for agent movement visualization
    \item \textbf{Connection Layer} (\texttt{connectionGraphics}): Visualizes inter-agent communication with animated arrows and particles
\end{itemize}

The shared canvas context (\texttt{ctx}) is implemented as a reactive object that maintains references to all PixiJS containers, sprite maps, and configuration state. This singleton pattern enables cross-composable access without prop drilling, following the Vue 3 Composition API best practices.

\subsubsection{Contagion Engine Composable Architecture}
\label{sec:contagion-engine}

The \texttt{useContagionEngine.js} composable encapsulates all simulation logic as a single, self-contained module that hooks into the existing render loop. The architecture follows the single responsibility principle, delegating rendering to existing systems while focusing exclusively on simulation state transitions.

\begin{figure}[t]
\centering
\begin{tikzpicture}[
    node distance=1.2cm and 1.5cm,
    box/.style={rectangle, draw, rounded corners, minimum width=2.8cm, minimum height=0.9cm, align=center, font=\small},
    data/.style={rectangle, draw, rounded corners, fill=blue!10, minimum width=2.2cm, minimum height=0.7cm, align=center, font=\footnotesize},
    arrow/.style={->, >=stealth, thick},
    darrow/.style={->, >=stealth, thick, dashed}
]
    % Main components
    \node[box, fill=green!15] (spatialcanvas) {SpatialCanvas.vue\\(Orchestrator)};
    \node[box, fill=orange!15, below=of spatialcanvas] (animationloop) {useAnimationLoop\\(Render Loop)};
    \node[box, fill=red!15, right=of animationloop] (contagion) {useContagionEngine\\(Simulation)};
    \node[box, fill=blue!15, left=of animationloop] (spatiallayout) {useSpatialLayout\\(Agent State)};
    \node[box, fill=purple!15, below=of animationloop] (spatialconfig) {useSpatialConfig\\(Configuration)};
    
    % Data stores
    \node[data, right=of contagion] (agentpos) {agentPositions\\(Map)};
    \node[data, right=of agentpos] (floors) {floorTiles\\(Array)};
    \node[data, below=of spatialconfig] (jsonconfig) {simulation\\\{JSON\}};
    
    % Connections
    \draw[arrow] (spatialcanvas) -- (animationloop);
    \draw[arrow] (animationloop) -- node[above, font=\scriptsize] {deltaMs} (contagion);
    \draw[arrow] (animationloop) -- (spatiallayout);
    \draw[arrow] (contagion) -- (agentpos);
    \draw[arrow] (contagion) -- (floors);
    \draw[arrow] (contagion) -- (spatialconfig);
    \draw[arrow] (spatialconfig) -- (jsonconfig);
    \draw[darrow] (spatiallayout) -- (agentpos);
    
    % Status system
    \node[data, below=of contagion] (status) {AGENT\_STATUS\\(Enum)};
    \draw[arrow] (contagion) -- (status);
    \draw[arrow] (spatiallayout) -- (status);
    
\end{tikzpicture}
\caption{System architecture diagram showing the integration of useContagionEngine with existing ChatDev spatial visualization components. Solid arrows indicate data flow; dashed arrows indicate read access.}
\label{fig:architecture}
\end{figure}

The contagion engine maintains singleton reactive state for:
\begin{itemize}
    \item \texttt{sandboxMode}: Boolean toggle for simulation activation
    \item \texttt{simulationRunning/simulationPaused}: Playback control state
    \item \texttt{elapsedTimeMs}: Cumulative simulation time
    \item \texttt{infectionTimers}: Map tracking infection timestamps per agent
    \item \texttt{immunityTimers}: Map tracking recovery timestamps for immunity expiry
    \item \texttt{floorDecayTimers}: Map tracking contamination decay timestamps per floor tile
\end{itemize}

The \texttt{updateContagion(deltaMs)} function is invoked once per frame within the render loop, following the established pattern of \texttt{updateIdleWanders()} and \texttt{updateAnimations()}. This integration point is illustrated in Figure~\ref{fig:architecture}.

\subsubsection{Agent Status State Machine}
\label{sec:status-machine}

The contagion simulation extends the existing \texttt{AGENT\_STATUS} enum with four epidemiological states:

\begin{enumerate}
    \item \texttt{HEALTHY}: Susceptible agents with green status glow (no pulse)
    \item \texttt{INFECTED}: Active infection with red status glow (rapid 3.0x pulse)
    \item \texttt{RECOVERED}: Post-infection immunity with blue status glow (gentle 0.5x pulse)
    \item \texttt{DECEASED}: Terminal state with gray status glow (no pulse, movement disabled)
\end{enumerate}

The state transitions follow a deterministic finite automaton:

\begin{equation}
\text{HEALTHY} \xrightarrow{\text{contact/proximity}} \text{INFECTED} \xrightarrow{\text{recovery}} \text{RECOVERED} \xrightarrow{\text{immunity expiry}} \text{HEALTHY}
\end{equation}

\begin{equation}
\text{INFECTED} \xrightarrow{\text{fatality roll}} \text{DECEASED}
\end{equation}

Each status maps to visual properties via the existing \texttt{STATUS\_COLORS} and \texttt{STATUS\_PULSE} dictionaries, ensuring consistent rendering through \texttt{useAgentRenderer.js}. The \texttt{EMOTE\_RULES} system triggers contextual emoji badges (e.g., ``Infected!'' with \texttt{:face\_with\_thermometer:}) based on state transitions.

\subsubsection{Config-Driven Simulation Parameters}
\label{sec:config-parameters}

All simulation parameters are externalized to the \texttt{SpatialConfig} JSON schema under the optional \texttt{simulation} key, enabling scenario customization without code changes. The schema definition is:

\begin{table}[h]
\centering
\caption{Simulation Configuration Parameters}
\label{tab:sim-params}
\begin{tabular}{@{}llll@{}}
\toprule
Parameter & Type & Default & Description \\
\midrule
\texttt{infectionRadius} & number & 80 & Proximity radius (px) \\
\texttt{infectionProbability} & number & 0.3 & Per-second transmission rate \\
\texttt{floorInfectionProbability} & number & 0.15 & Floor-to-agent base rate \\
\texttt{recoveryTimeMs} & number & 15000 & Infection duration (ms) \\
\texttt{fatalityProbability} & number & 0.05 & Death probability on recovery \\
\texttt{contaminationDecayMs} & number & 10000 & Decay interval (ms) \\
\texttt{immunityDurationMs} & number & 30000 & Immunity duration (ms) \\
\bottomrule
\end{tabular}
\end{table}

The \texttt{FloorTile} type extends with an optional \texttt{contaminationLevel} field (0--3), driving color overlay gradients from green (clean) through yellow and orange to red (severe). Example configuration:

\begin{verbatim}
{
  "simulation": {
    "infectionRadius": 80,
    "infectionProbability": 0.3,
    "recoveryTimeMs": 15000
  },
  "floors": [{
    "id": "contaminated_zone",
    "position": {"x": 200, "y": 200},
    "width": 120,
    "height": 120,
    "contaminationLevel": 3
  }]
}
\end{verbatim}

Parameters are loaded at initialization via \texttt{useSpatialConfig.js}, with fallback to hardcoded defaults for backward compatibility with configurations lacking the \texttt{simulation} block.

\subsubsection{Frame-Rate Independence}
\label{sec:frame-rate}

To ensure consistent simulation behavior across varying frame rates, all probabilistic events use time-scaled calculations. The probability scaling formula is:

\begin{equation}
P_{\text{scaled}} = 1 - (1 - P_{\text{base}})^{\Delta t / 1000}
\end{equation}

where $P_{\text{base}}$ is the per-second probability from configuration and $\Delta t$ is the frame delta in milliseconds. This approach maintains expected infection rates regardless of frame timing jitter.

\subsubsection{Proximity and Contact Detection}
\label{sec:proximity-detection}

The simulation employs O($n^2$) pairwise proximity checks for agent-to-agent transmission, consistent with the existing \texttt{applyPerFrameSeparation()} function used for collision avoidance. For $n < 20$ agents---the typical range for ChatDev multi-agent workflows---this approach incurs negligible overhead (400 distance calculations per frame at 60 FPS = 24,000 operations/second).

Floor contact detection uses axis-aligned bounding box (AABB) intersection tests between agent positions and floor tile rectangles, with contamination probability scaled by level: $P_{\text{floor}} = P_{\text{base}} \times (\text{level} / 3)$.

\subsubsection{Integration with Existing Render Loop}
\label{sec:render-loop-integration}

The \texttt{useAnimationLoop.js} composable orchestrates per-frame updates through a requestAnimationFrame (RAF) callback chain. The contagion simulation integrates as a single additional call:

\begin{verbatim}
function renderLoop() {
  updateActiveStates()
  updateIdleWanders()
  if (updateContagion) updateContagion(deltaMs)
  updateAnimations()
  applyPerFrameSeparation()
  updateEmotes()
  drawTrailParticles()
  drawConnections()
  cleanupConnections()
  cleanupTrailParticles()
}
\end{verbatim}

This design ensures that simulation updates occur in lockstep with rendering, preventing desynchronization between agent visual positions and simulation state. The sandbox mode toggle gates all simulation behavior, ensuring zero performance impact when the contagion feature is inactive.
